\chapter{Eigenvalues and Invariant Subspaces}\label{linear:ch:eigen}

\begin{center}\itshape
I have not thought it necessary to undertake the labour of a formal proof of the theorem in the general case of a
matrix of any degree.
\upshape(Arthur Cayley, \emph{A Memoir on the Theory of Matrices}, 1858)
\end{center}
\noindent\emph{Use Volume~1 and Appendices~\ref{linear:ch:duality} and~\ref{linear:ch:volume}. This gives an algebraic companion to Chapter~\ref{linear:ch:spectral} and supports the dynamics in Appendix~\ref{linear:ch:dynamics}; Volume~9 uses the matrix exponential.}

Planetary orbits slowly change because the planets perturb one another. Lagrange and Laplace studied these secular variations in the eighteenth century by linearising the equations. Their slow oscillation frequencies came from roots of a secular equation. Cauchy's 1829 work proved reality of the roots for symmetric systems, an early spectral theorem. That linearised stability question was a major motivation, though it did not settle every nonlinear question about the solar system.

Sylvester named the matrix in 1850. Cayley's 1858 memoir treated it as an object with its own operations and asserted its characteristic equation. He checked small sizes but declined a general proof in the words above; Frobenius supplied one in 1878. Weierstrass and Jordan developed normal forms, and Jordan's 1874 dispute with Kronecker concerned what a canonical form should accomplish. Our question is concrete: what does a map look like after the best possible basis change?

Invariant subspaces answer it in stages. A complete eigenbasis diagonalises the map. When that fails, triangularisation still works over a splitting field, and generalised eigenspaces separate the eigenvalues. On each piece, what remains is nilpotent. The dimensions of its successive kernels determine Jordan chains, with their lengths recorded by a Young diagram.

\begin{think}
Classify $4\times4$ complex matrices satisfying $N^4=0$ up to similarity. Then impose $N^2=0$. The first problem has five answers, one per partition of four; the second keeps the three partitions with parts at most two. Try using the ranks of $N,N^2,N^3$ to distinguish them before reading Section~\ref{linear:sec:nilpotent}.
\end{think}

\section{Invariant subspaces}

Let $A$ act on a finite-dimensional space $V$ over $\mathbb K$. A subspace $W$ is \emph{invariant} if $A(W)\subseteq W$. We then have both the restriction $A|_W$ and the quotient operator $\bar A(v+W)=Av+W$. Invariance makes the quotient formula well defined. Extending a basis of $W$ gives a block matrix $\begin{pmatrix}A|_W&*\\0&\bar A\end{pmatrix}$, so
\begin{equation}\label{linear:eq:charfactor}
\chi_A=\chi_{A|_W}\chi_{\bar A},\qquad\chi_A(X)=\det(XI-A).
\end{equation}
A nonzero vector with $Av=\lambda v$ is an eigenvector, and its line is invariant. Its eigenvalue is a root of $\chi_A$. Kernels and images of polynomials in $A$ are invariant because every $P(A)$ commutes with $A$. The zero-dimensional case has characteristic polynomial one and an empty diagonal or Jordan form; assertions requiring a nonzero eigenvector will explicitly exclude it.

\begin{proposition}\label{linear:prop:indep}
Eigenvectors for distinct eigenvalues are independent. In particular, $n$ distinct characteristic roots in $\mathbb K$ give a diagonalisation in dimension $n$.
\end{proposition}
\begin{proof}
Induct on the number $k$ of vectors. Apply $A-\lambda_k I$ to a relation $\sum_ic_iv_i=0$. Independence of the first $k-1$ eigenvectors forces $c_i(\lambda_i-\lambda_k)=0$ for $i<k$, hence those coefficients vanish. The remaining nonzero vector forces $c_k=0$.
\end{proof}

\section{Triangularisation}

A polynomial \emph{splits} over $\mathbb K$ when it is a product of linear factors there. Volume~1, Chapter~6, gives the fundamental theorem of algebra, so complex polynomials split.

\begin{theorem}[triangularisation]\label{linear:thm:schur}
An operator admits an upper triangular matrix exactly when its characteristic polynomial splits over its field. Thus every complex matrix is triangularisable. Over $\R$ or $\C$ with an inner product, a split characteristic polynomial permits an orthonormal triangularising basis, Schur's theorem.
\end{theorem}
\begin{proof}
A triangular characteristic polynomial is the product of its diagonal linear factors. Conversely, induct on dimension. A nonzero split space has an eigenvector $v_1$. Equation~\eqref{linear:eq:charfactor} makes the induced characteristic polynomial on $V/\mathbb Kv_1$ split too. Triangularise that quotient by induction and lift its basis. Together with $v_1$ the lifts form a basis, whose successive spans are invariant, giving an upper triangular matrix.

For the orthonormal version, use the inner-product projection already taught in Volume~1 and main Chapter~\ref{ch:hilbert}. Take $v_1$ unit and compress $A$ to its orthogonal complement. In a basis beginning with $v_1$, the lower block is this compression and its characteristic polynomial is $\chi_A/(X-\lambda)$. Induct on it and adjoin $v_1$. No result from the later inner-product appendix is needed for this step.
\end{proof}

The diagonal entries list eigenvalues with algebraic multiplicity, giving trace as their sum and determinant as their product. A polynomial of a triangular matrix has diagonal $P(\lambda_i)$, proving polynomial spectral mapping as well. Over $\R$, a quarter-turn has no invariant line and cannot be triangularised.

\section{The Cayley--Hamilton theorem}
\begin{theorem}[Cayley--Hamilton]\label{linear:thm:cayleyhamilton}
Every square matrix over a commutative ring $R$ satisfies $\chi_A(A)=0$.
\end{theorem}
\begin{proof}
The positive-size case has $B(X)=\operatorname{adj}(XI-A)=\sum_{k=0}^{n-1}B_kX^k$. The adjugate identity over $R[X]$ gives $(XI-A)B(X)=\chi_A(X)I$. If $\chi_A=\sum_{k=0}^nc_kX^k$, comparison yields $B_{k-1}-AB_k=c_kI$ for $0\le k\le n$, with $B_{-1}=B_n=0$. Multiply on the left by $A^k$ and sum. The left side telescopes:
\[\sum_{k=0}^n(A^kB_{k-1}-A^{k+1}B_k)=0=\sum_{k=0}^nc_kA^k.
\]
The empty-matrix case is immediate, since its identity map is the zero map.
\end{proof}

Substituting $X=A$ directly into the scalar determinant $\det(XI-A)$ would confuse a matrix polynomial with a scalar determinant. The coefficient proof specifies the order of multiplication instead.

Over a field, annihilating polynomials form a nonzero ideal. Its monic generator, the \emph{minimal polynomial} $\mu_A$, divides every annihilator by Euclidean division: a smaller nonzero remainder cannot also kill $A$. In particular $\mu_A\mid\chi_A$. Every eigenvalue is a root of $\mu_A$, since applying it to its eigenvector gives $\mu_A(\lambda)v=0$. The two polynomials thus have the same roots in a splitting field, with possibly different multiplicities.

\section{Generalised eigenspaces}

The matrix $\begin{pmatrix}0&1\\0&0\end{pmatrix}$ has too few eigenvectors for diagonalisation. Kernels of higher powers enlarge eigenspaces enough to decompose the whole space.

\begin{lemma}[kernel lemma]\label{linear:lem:kernels}
For pairwise coprime polynomials $P_i$ and $P=\prod_iP_i$,
\[\ker P(A)=\bigoplus_i\ker P_i(A).\]
The summand projections on this kernel are restrictions of polynomials in $A$.
\end{lemma}
\begin{proof}
If a factor is zero, pairwise coprimality makes every other factor a nonzero constant; the decomposition and projections are then $V$ and zero summands, with polynomial projections $I,0$. Otherwise put $Q_i=\prod_{j\ne i}P_j$. Their common gcd is one, so B\'ezout gives $\sum_iU_iQ_i=1$. Define $E_i=U_i(A)Q_i(A)$. For $v\in\ker P(A)$, $v=\sum_iE_iv$ and $P_i(A)E_iv=0$, proving existence of the decomposition.

On $\ker P_i(A)$, every $E_j$ with $j\ne i$ vanishes because $Q_j$ contains $P_i$. The sum identity then makes $E_i$ the identity there. Applying $E_i$ to any sum of these kernel vectors proves uniqueness. These are exactly the claimed projections. The empty product case has kernel zero and the empty direct sum.
\end{proof}

\begin{theorem}[generalised eigenspaces]\label{linear:thm:geneigen}
Suppose $\chi_A=\prod_i(X-\lambda_i)^{m_i}$ with distinct $\lambda_i$. Set $V_i=\ker(A-\lambda_i I)^{m_i}$. Then $V=\bigoplus_iV_i$, every $V_i$ is invariant with dimension $m_i$, and $A-\lambda_i I$ is nilpotent on it. Also $V_i=\ker(A-\lambda_i I)^{n'}$ for every $n'\ge m_i$.
\end{theorem}
\begin{proof}
Cayley--Hamilton and the kernel lemma give the invariant direct sum, with the stated nilpotence. Each restricted characteristic polynomial divides $\chi_A$ by the invariant block factorisation, so it splits. Its only possible eigenvalue is $\lambda_i$, since a different eigenvector value would contradict $(A-\lambda_i I)^{m_i}=0$. Thus its characteristic polynomial is $(X-\lambda_i)^{d_i}$, where $d_i=\dim V_i$. Multiplying the block polynomials gives $d_i=m_i$.

On $V_j$ with $j\ne i$, $A-\lambda_i I$ is invertible, either by its nonzero determinant or a finite inverse for scalar plus nilpotent. Hence the kernel of any of its powers lies in $V_i$. For powers at least $m_i$ it contains all of $V_i$, proving stabilisation.
\end{proof}

\begin{corollary}[diagonalisability]\label{linear:cor:diag}
The following are equivalent: $A$ is diagonalisable over $\mathbb K$; its minimal polynomial splits with simple roots; some split polynomial with simple roots annihilates it.
\end{corollary}
\begin{proof}
A diagonal map is killed by the product of its distinct eigenvalue factors. Conversely, applying the kernel lemma to a simple-root annihilator decomposes $V$ into eigenspaces. The minimal polynomial divides every annihilator, making the middle condition equivalent to the last.
\end{proof}

\begin{theorem}[Dunford--Jordan--Chevalley decomposition]\label{linear:thm:dunford}
If $\chi_A$ splits, there is exactly one decomposition $A=D+N$ with $D$ diagonalisable, $N$ nilpotent and $DN=ND$. Both parts are polynomials in $A$.
\end{theorem}
\begin{proof}
Use the polynomial projections $E_i$ onto $V_i$ to define $D=\sum_i\lambda_iE_i$ and $N=A-D$. They are polynomials in $A$, commute, and have the desired restrictions on each summand, proving existence.

For another pair $D',N'$, the operator $D'$ commutes with $A$, hence with $D$ and $N$. Two commuting diagonalisable maps have a common eigenbasis: each eigenspace of the first is preserved by the second, whose restriction is still killed by its simple-root minimal polynomial and hence diagonalises by the preceding corollary. This proves the needed simultaneous-diagonalisation fact here. Therefore $D-D'$ is diagonalisable. Also $N'$ commutes with $N$, and, choosing $q$ with $N^q=N'^q=0$, the binomial expansion gives $(N'-N)^{2q}=0$. Since $D-D'=N'-N$ is both diagonal and nilpotent, it is zero. The zero-space case is immediate. Thus both parts are unique.
\end{proof}

\section{Nilpotent operators and Young diagrams}\label{linear:sec:nilpotent}

On a generalised eigenspace, subtracting its eigenvalue leaves a nilpotent operator. For a nonzero space let $s$ be its nilpotence index, the least positive integer with $N^s=0$. The kernels $K_k=\ker N^k$ grow strictly until $K_s=V$. Indeed, equality at one stage would imply equality at the next by applying $N$, and then forever, contradicting the minimal index. Put $r_k=\dim K_k-\dim K_{k-1}$.

\begin{lemma}\label{linear:lem:decreasing}
The induced maps $K_{k+1}/K_k\to K_k/K_{k-1}$, given by $N$, are injective. Consequently $r_1\ge\cdots\ge r_s>0$ and $\sum_{k=1}^sr_k=\dim V$ for nonzero $V$. For the zero space the list and sum are empty.
\end{lemma}
\begin{proof}
The indicated quotient map is well defined. If $Nv\in K_{k-1}$, then $N^kv=0$, so $v\in K_k$; its kernel is zero. Quotient dimensions give the inequalities, and their sum telescopes.
\end{proof}

Draw columns with heights $r_1,r_2,\ldots$. Their non-increase makes a Young diagram; its row lengths form the conjugate partition. Figure~\ref{linear:fig:young} gives an example. The nilpotent block $J_m$ has ones above its diagonal, sending $e_1$ to zero and $e_j$ to $e_{j-1}$ for $j>1$.
\begin{figure}[ht]
\centering
\fig{linear-ch03-young}
\caption{The Young diagram of a nilpotent operator with $\dim\ker N^k=4,7,8,9$ for $k=1,2,3,4$. Each row is a Jordan
chain $u,Nu,N^2u,\dots$ read from right to left, ending in the kernel of $N$ (first column); here the Jordan blocks have
sizes $4,2,2,1$.}
\label{linear:fig:young}
\end{figure}

\begin{theorem}[Jordan form of a nilpotent operator]\label{linear:thm:jordannil}
A nilpotent $N$ has a basis in which it is $\operatorname{diag}(J_{m_1},\ldots,J_{m_t})$, with decreasing block sizes equal to the Young diagram's row lengths. The number of blocks of size at least $k$ is $r_k$, so these sizes are uniquely determined. On the zero space the block list is empty.
\end{theorem}
\begin{proof}
Choose chain tops from the highest level down. At level $k$, the classes of the vectors $N^{j-k}u_{j,i}$ from already selected tops at levels $j>k$ are independent in $K_k/K_{k-1}$. They are images of independent classes from the next quotient under the preceding injective map. Complete them to a quotient basis with new tops $u_{k,i}\in K_k$.

Take every chain $N^{k-1}u_{k,i},\ldots,Nu_{k,i},u_{k,i}$. At each kernel level their classes form its quotient basis by construction. The full collection is independent: in a relation, passing to the highest level present forces all coefficients there to vanish, then descend. Its size is $\sum r_k=\dim V$, so it is a basis. Each ordered chain has matrix $J_k$, proving existence.

For any such block list, $\dim\ker N^k=\sum_i\min(k,m_i)$. Subtract consecutive dimensions to see that $r_k$ counts exactly the blocks with $m_i\ge k$. Hence the kernel dimensions determine the list, proving uniqueness.
\end{proof}

\begin{corollary}[Jordan normal form]\label{linear:cor:jordan}
If $\chi_A$ splits, $A$ is similar to a block diagonal matrix with blocks $\lambda I_m+J_m$, unique up to block order. The number of blocks for $\lambda$ of size at least $k$ is
\[\operatorname{rank}(A-\lambda I)^{k-1}-\operatorname{rank}(A-\lambda I)^k.
\]
Two such matrices are similar exactly when their eigenvalues and all these power ranks agree.
\end{corollary}
\begin{proof}
Apply the nilpotent theorem on each generalised eigenspace. On every other eigenvalue summand, $A-\lambda I$ is invertible, so its power kernels come only from the $\lambda$ summand. Rank--nullity gives the formula. Similarity preserves those ranks; conversely they recover all block sizes, hence the same Jordan form.
\end{proof}

\begin{corollary}\label{linear:cor:transpose}
A complex square matrix and its transpose belong to the same similarity class.
\end{corollary}
\begin{proof}
The characteristic polynomials agree. Also $(A^{\mathsf T}-\lambda I)^k=((A-\lambda I)^k)^{\mathsf T}$ has the same rank by Appendix~\ref{linear:ch:duality}. The Jordan criterion gives similarity.
\end{proof}

\begin{remark}
When the characteristic polynomial does not split, modules over $\mathbb K[X]$ give rational canonical form instead. Volume~3, Chapter~7, develops that route and proves transpose similarity over every field. Here the kernel-quotient construction directly explains the Jordan chains: column heights are kernel increments and row lengths are chain lengths.
\end{remark}

Jordan structure can change abruptly under small perturbations. Adding $\varepsilon$ in the bottom-left corner of $J_n$ gives characteristic polynomial $X^n-\varepsilon$, so the eigenvalues lie on a circle of radius $|\varepsilon|^{1/n}$. For $n=12,\varepsilon=10^{-6}$, this radius is exactly $10^{-1/2}>0.3$. Figure~\ref{linear:fig:perturb} shows the effect. This sensitivity makes exact Jordan classification unsuitable for routine floating-point computation; Schur and singular-value methods provide better numerical representations. Stability of a particular algorithm still requires its own error analysis.
\begin{figure}[ht]
\centering
\fig{linear-ch03-perturb}
\caption{The eigenvalues of $J_{12}+10^{-6}E_{12,1}$: a perturbation of size $10^{-6}$ moves the twelve-fold eigenvalue $0$
to a circle of radius $10^{-1/2}\approx0.32$.}
\label{linear:fig:perturb}
\end{figure}

\section{Commuting operators}
\begin{theorem}\label{linear:thm:commuting}
For a pairwise commuting family $\mathcal F$ on finite-dimensional $V$:
\begin{enumerate}[label=(\alph*)]
\item if every member is diagonalisable, there is a common diagonalising basis;
\item if every characteristic polynomial splits, there is a common triangularising basis and, when $V\ne\{0\}$, a common eigenvector.
\end{enumerate}
\end{theorem}
\begin{proof}
(a) Induct on dimension. If every member is scalar, any basis works. Otherwise choose a nonscalar diagonal member $A$ and split into its eigenspaces. Each has smaller dimension and is preserved by the family, since commutation preserves eigenvectors' eigenvalues. Restrictions of a diagonal member remain killed by its simple-root annihilator, so are diagonalisable. Induct on each summand and join the bases. This also proves the assertion for an infinite family; the induction is on dimension, not family size.

(b) On a nonzero space choose a nonzero common invariant subspace $W$ of least dimension. Each restricted characteristic polynomial divides the split polynomial on $V$, so has an eigenvalue. If any restriction were nonscalar, one of its proper nonzero eigenspaces would be invariant under the whole commuting family, contradicting minimality. Thus every member acts as a scalar on $W$, and any nonzero vector there is common. Quotient by its line; the induced family still commutes and has split characteristic polynomials. Dimension induction gives simultaneous triangularisation by lifting the quotient basis. On the zero space the empty basis triangularises the family, but no eigenvector exists.
\end{proof}

\section{Functions of a matrix}

For a polynomial over any field, write $c_j(\lambda)=[X^j]f(\lambda+X)$. On a Jordan block, nilpotence gives
\begin{equation}\label{linear:eq:fjordan}
\begin{aligned}
f(\lambda I+J_m)&=\sum_{j=0}^{m-1}c_j(\lambda)J_m^j\\
&=\begin{pmatrix}c_0&c_1&c_2&\cdots\\&c_0&c_1&\ddots\\&&\ddots&\ddots\\&&&c_0\end{pmatrix},\qquad
c_j(\lambda)=\frac{f^{(j)}(\lambda)}{j!}\ \text{in characteristic zero}.
\end{aligned}
\end{equation}
The coefficient definition works even when a factorial vanishes in the field. It shows that $f(A)$ depends only on the polynomial remainder modulo the powers $(X-\lambda)^{s_\lambda}$, where $s_\lambda$ is the largest block size. In particular those powers give the minimal-polynomial exponents, also explored in Exercise~\ref{linear:ex:minpolyjordan}.

Over $\C$, a function holomorphic near the eigenvalues has the same finite Taylor jets and defines $f(A)$ blockwise. To see basis independence, the kernel lemma or polynomial Chinese remainders give one polynomial with those jets, and all choices differ by a multiple of the minimal polynomial. This definition agrees with a convergent power series by norm convergence and the finite block expansion. For real eigenvalues, a sufficiently smooth real function can likewise supply real derivatives. Mere real smoothness on a complex neighbourhood does not supply complex derivatives; holomorphy is the condition used here.

\begin{proposition}[the matrix exponential]\label{linear:prop:exp}
For complex square $A$, the series $e^A=\sum_{k\ge0}A^k/k!$ converges absolutely in each entry. It satisfies
\[e^{A+B}=e^Ae^B\ (AB=BA),\qquad e^{PAP^{-1}}=Pe^AP^{-1},\]
\[\det e^A=e^{\operatorname{tr}A},\qquad\frac{\dd}{\dd t}e^{tA}=Ae^{tA}.
\]
It is invertible with inverse $e^{-A}$.
\end{proposition}
\begin{proof}
For positive matrix size, the norm $\|A\|=n\max_{i,j}|a_{ij}|$ is submultiplicative. The exponential series is therefore dominated by the scalar series $\sum\|A\|^k/k!$. For commuting matrices the binomial formula and absolutely convergent Cauchy product prove the addition rule, hence the inverse rule. Conjugation passes through partial sums and their limit.

Triangularising $A$ makes $e^A$ triangular with diagonal $e^{\lambda_i}$; their product is $e^{\sum\lambda_i}=e^{\operatorname{tr}A}$. Finally the addition rule and $(e^{hA}-I)/h\to A$ prove the derivative. The remainder here is norm bounded by $|h|\|A\|^2e^{|h|\|A\|}/2$. All assertions hold for the empty matrix by their empty-sum and determinant conventions.
\end{proof}

A skew-symmetric two-by-two generator exponentiates to a rotation (Exercise~\ref{linear:ex:rotation}). Appendix~\ref{linear:ch:dynamics} uses $e^{tA}$ for differential equations, and Volume~9 develops the associated Lie-theoretic viewpoint.

\begin{beauty}
One kernel dimension gives only a rough measurement. The whole sequence of kernel dimensions, however, reconstructs every nilpotent chain length. This converts similarity classification into partitions and makes the Young diagram an exact invariant, rather than a picture attached after the calculation. Related diagrams index symmetric-group representations in Volume~10.
\end{beauty}

\begin{figure}[ht]
\centering\fig{appC-rw-square}
\caption{Squaring a length-four nilpotent chain separates its odd and even positions into two length-two chains. In general a block of size $m$ splits under squaring into lengths $\lceil m/2\rceil$ and $\lfloor m/2\rfloor$, omitting a zero length.}
\label{fig:v8appCsquare}
\end{figure}
\begin{figure}[ht]
\centering\fig{appC-rw-parts}
\caption{On each generalised eigenspace, the diagonal part is its scalar eigenvalue and the remaining part is nilpotent. Polynomial projections assemble these local pieces into $A=D+N$ with $DN=ND$.}
\label{fig:v8appCparts}
\end{figure}
\section{Exercises}
\begin{exercise}[$\star$]\label{linear:ex:abba}
Prove $\chi_{AB}=\chi_{BA}$ for square $A,B$ over any field, starting with invertible $A$. For rectangular $A\in M_{m,n}$, $B\in M_{n,m}$ with $m\le n$, establish $\chi_{BA}=X^{n-m}\chi_{AB}$.
\end{exercise}
\begin{exercise}[$\star$]\label{linear:ex:companion}
Let $P=X^n+a_{n-1}X^{n-1}+\cdots+a_0$. Its companion has ones below the diagonal and last column $(-a_0,\ldots,-a_{n-1})^{\mathsf T}$. Prove that its characteristic and minimal polynomials both equal $P$.
\end{exercise}
\begin{exercise}[$\star$]\label{linear:ex:nilcount}
Count complex similarity classes with $N^2=0$ in dimension $n$. Then count the classes with $N^3=0$ in dimension six.
\end{exercise}
\begin{exercise}[$\star\star$]\label{linear:ex:jordanfind}
A complex $7\times7$ matrix has $\chi_A=(X-2)^5(X+1)^2$, ranks of $(A-2I)^k$ equal to $4,3,2$ for $k=1,2,3$, and $\operatorname{rank}(A+I)=6$. Determine its Jordan form and minimal polynomial.
\end{exercise}
\begin{exercise}[$\star\star$]\label{linear:ex:finiteorder}
If complex $A$ satisfies $A^k=I$ for some positive integer $k$, prove it diagonalisable. Test the claim over $\mathbb F_2$ and explain the difference.
\end{exercise}
\begin{exercise}[$\star\star$]\label{linear:ex:nilptrace}
In characteristic zero, prove nilpotence equivalent to $\operatorname{tr}(A^k)=0$ for every $1\le k\le n$.
\end{exercise}
\begin{exercise}[$\star\star$]\label{linear:ex:sqrt}
Show that the block $J_n$, $n\ge2$, has no complex square root, whereas every $\lambda I+J_n$ with $\lambda\ne0$ has one.
\end{exercise}
\begin{exercise}[$\star\star$]\label{linear:ex:commutant}
For complex $A$, show its commuting matrices form a subspace of dimension at least $n$. Prove equality when the eigenvalues are distinct, and identify that commutant as the polynomials in $A$.
\end{exercise}
\begin{exercise}[$\star$]\label{linear:ex:rotation}
Evaluate $\exp\begin{pmatrix}0&-t\\t&0\end{pmatrix}$ and $\exp(t(\lambda I+J_3))$ exactly.
\end{exercise}
\begin{exercise}[$\star\star$]\label{linear:ex:minpolyjordan}
For split $\chi_A$, prove $\mu_A=\prod_\lambda(X-\lambda)^{s_\lambda}$, with $s_\lambda$ the largest block size at $\lambda$. Deduce that $\mu_A=\chi_A$ precisely when every eigenvalue has a single Jordan block. The proof must work in positive characteristic as well.
\end{exercise}
\begin{exercise}[$\star\star$]\label{linear:ex:expsum}
Give two-by-two examples where $e^{A+B}\ne e^Ae^B$, and a nonzero matrix with $e^A=I$.
\end{exercise}
\begin{exercise}[$\star\star\star$]\label{linear:ex:symprod}
(a) For the reversal matrix $R$, prove $RJ_mR=J_m^{\mathsf T}$. (b) Construct, for every complex square $A$, an invertible symmetric $S$ with $A^{\mathsf T}=S^{-1}AS$. (c) Factor $A$ as two symmetric matrices, one invertible, as in Frobenius's 1910 theorem.
\end{exercise}

\begin{exercise}[$\star\star$]\label{ex:v8cCa}
\emph{Polynomial projections with a repeated root.} Suppose $\mu_A=(X-1)^2(X+1)$ over $\C$. Prove that $E_-=(A-I)^2/4$ and $E_+=I-E_-$ are the generalised-eigenspace projections. Compute $D=E_+-E_-$ and show $N=A-D$ satisfies $N^2=0$.
\end{exercise}
\begin{exercise}[$\star\star$]\label{ex:v8cCb}
\emph{Squaring the chains.} A nilpotent operator on dimension ten has kernel dimensions $4,7,9,10$ for its first four powers. Find its block sizes, then those of its square. Prove the rule splitting a size-$m$ block into lengths $\lceil m/2\rceil$ and $\lfloor m/2\rfloor$.
\end{exercise}
\begin{exercise}[$\star\star$]\label{ex:v8cCc}
\emph{One chain's commutant.} Over any field, prove every matrix commuting with $J_n$ is a polynomial in $J_n$, uniquely of degree below $n$. Give its upper triangular Toeplitz form and explain why its last column determines it.
\end{exercise}
\begin{exercise}[$\star\star$]\label{ex:v8cCd}
\emph{Exponentiation retains nilpotent chains.} For complex nilpotent $N$, prove $e^N-I=NR(N)$ with $R(0)=1$. Show that $R(N)$ is invertible and all power kernels of $e^N-I$ equal those of $N$. Deduce that $e^N$ has the same block sizes at eigenvalue one as $N$ has at zero.
\end{exercise}

\section{Hints}
\begin{hint}{linear:ex:abba}
Compare $\begin{pmatrix}AB&0\\B&0\end{pmatrix}$ and $\begin{pmatrix}0&0\\B&BA\end{pmatrix}$ by the shear $\begin{pmatrix}I&A\\0&I\end{pmatrix}$.
\end{hint}
\begin{hint}{linear:ex:companion}
The vectors $e_1,Ce_1,\ldots,C^{n-1}e_1$ are a basis, while $P(C)e_1=0$.
\end{hint}
\begin{hint}{linear:ex:nilcount}
Translate the power condition into a maximum allowed part size in a partition.
\end{hint}
\begin{hint}{linear:ex:jordanfind}
Subtract consecutive ranks to obtain the Young column heights, keeping the other eigenvalue's invertible block in mind.
\end{hint}
\begin{hint}{linear:ex:finiteorder}
Inspect the roots and derivative of $X^k-1$. A nontrivial unipotent two-by-two matrix gives a characteristic-two test.
\end{hint}
\begin{hint}{linear:ex:nilptrace}
Pass to a splitting field and apply a Vandermonde system to the distinct nonzero eigenvalues with their multiplicities.
\end{hint}
\begin{hint}{linear:ex:sqrt}
A square root would be nilpotent too. Compare its first two kernel dimensions; for the nonzero eigenvalue case truncate a square-root binomial series.
\end{hint}
\begin{hint}{linear:ex:commutant}
Independent polynomials on each Jordan block already give $n$ commuting matrices. For distinct eigenvalues use diagonal coordinates and interpolation.
\end{hint}
\begin{hint}{linear:ex:rotation}
Separate the even and odd powers of a matrix squaring to $-I$.
\end{hint}
\begin{hint}{linear:ex:minpolyjordan}
Expand $P(\lambda+X)$ in coefficients, rather than dividing derivatives by factorials in a field where they may vanish.
\end{hint}
\begin{hint}{linear:ex:expsum}
Try $E_{12},E_{21}$. Their sum's exponential has equal diagonal entries, while the product of their exponentials does not.
\end{hint}
\begin{hint}{linear:ex:symprod}
If $A=PJP^{-1}$ and $RJR=J^{\mathsf T}$, use the symmetric matrix $PRP^{\mathsf T}$.
\end{hint}
\begin{hint}{ex:v8cCa}
Reduce the proposed projection polynomials modulo $(X-1)^2$ and $X+1$ separately.
\end{hint}
\begin{hint}{ex:v8cCb}
Subtract kernel dimensions, then read rows from columns. Under $N^2$, odd and even chain positions separate.
\end{hint}
\begin{hint}{ex:v8cCc}
Since $e_n,J_ne_n,\ldots,J_n^{n-1}e_n$ is a basis, commutation propagates the image of $e_n$ to every vector.
\end{hint}
\begin{hint}{ex:v8cCd}
The exponential is a finite polynomial on a nilpotent. A polynomial with nonzero constant term is invertible there by a finite geometric series.
\end{hint}

\section{Solutions}
\begin{solution}{linear:ex:abba}
\idea A block shear gives the relation exactly, even when neither factor is invertible.
\writeup With invertible $A$, $AB=A(BA)A^{-1}$. In general set
\[M=\begin{pmatrix}AB&0\\B&0\end{pmatrix},\quad
M'=\begin{pmatrix}0&0\\B&BA\end{pmatrix},\quad P=\begin{pmatrix}I_m&A\\0&I_n\end{pmatrix}.
\]
Multiplication gives $MP=PM'$. The shear $P$ is invertible, so block determinants yield $X^n\chi_{AB}=X^m\chi_{BA}$. Comparing coefficients after cancelling $X^m$ gives the rectangular identity; at $m=n$ it gives equality. The computation uses no division and also works over commutative rings.
\end{solution}
\begin{solution}{linear:ex:companion}
\idea A cyclic vector forces every annihilator to have full degree.
\writeup The first $n$ powers of $C$ on $e_1$ give $e_1,\ldots,e_n$, so no nonzero polynomial of degree below $n$ can kill $C$. The last-column relation gives $P(C)e_1=0$. Since every basis vector is a power of $C$ on $e_1$, commutation makes $P(C)$ zero everywhere. Thus $\mu_C=P$. Cayley--Hamilton makes this divide the monic degree-$n$ characteristic polynomial, forcing $\chi_C=P$ as well.
\end{solution}
\begin{solution}{linear:ex:nilcount}
\idea The largest allowed block equals the allowed nilpotent exponent.
\writeup For $N^2=0$, a partition has only ones and twos. Choosing the number $j$ of twos, $0\le j\le\lfloor n/2\rfloor$, determines it, giving $\lfloor n/2\rfloor+1$ classes. For dimension six with $N^3=0$, the partitions are $3+3$, $3+2+1$, $3+1+1+1$, $2+2+2$, $2+2+1+1$, $2+1+1+1+1$, and six ones. Thus there are seven classes.
\end{solution}
\begin{solution}{linear:ex:jordanfind}
\idea Rank drops count blocks that reach each chain level.
\writeup At eigenvalue two, ranks $7,4,3,2$ give column heights $3,1,1$. Their sum is five, exhausting that eigenvalue's multiplicity. The row lengths are therefore $3,1,1$. At eigenvalue minus one, rank six means one kernel vector and hence one block, of length two. Thus
\[A\sim\operatorname{diag}(2I_3+J_3,2,2,-I_2+J_2).
\]
The displayed blocks are killed exactly by multiples of $(X-2)^3(X+1)^2$, so that is $\mu_A$. This conclusion follows directly from block substitution and does not require a later exercise.
\end{solution}
\begin{solution}{linear:ex:finiteorder}
\idea A simple-root annihilator excludes every nontrivial Jordan extension.
\writeup The polynomial $X^k-1$ has no common root with its derivative $kX^{k-1}$ over $\C$. It splits with simple roots and kills $A$, so Corollary~\ref{linear:cor:diag} applies. Over $\mathbb F_2$, $A=\begin{pmatrix}1&1\\0&1\end{pmatrix}$ satisfies $A^2=I$ but is not diagonal: its sole eigenvalue is one, and diagonalisation would make it the identity. Here $X^2-1=(X-1)^2$, so the simple-root argument fails.
\end{solution}
\begin{solution}{linear:ex:nilptrace}
\idea Distinct nonzero eigenvalues form an invertible Vandermonde system whose weights are their multiplicities.
\writeup Pass to an algebraic splitting field, as constructed in Volume~3, without changing the traces or matrix identities. Let the nonzero distinct eigenvalues be $\mu_1,\ldots,\mu_r$ with positive integer multiplicities $n_j$. Triangularisation gives $\operatorname{tr}(A^k)=\sum_jn_j\mu_j^k$. Vanishing for $k=1,\ldots,r$ makes the vector $(n_j\mu_j)_j$ lie in the kernel of the Vandermonde matrix $(\mu_j^{k-1})_{k,j}$. It is invertible, forcing $n_j\mu_j=0$. Characteristic zero makes every $n_j$ nonzero in the field, a contradiction unless $r=0$. Thus $\chi_A=X^n$ and Cayley--Hamilton gives $A^n=0$. Conversely a nilpotent matrix triangularises with zero diagonal, as do its positive powers, giving all traces zero. Characteristic is essential: $I_p$ in characteristic $p$ has every power trace zero but is not nilpotent.
\end{solution}
\begin{solution}{linear:ex:sqrt}
\idea A square root would have a kernel growing too slowly for a nonzero nilpotent operator.
\writeup If $B^2=J_n$, then $B^{2n}=0$, so $B$ is nilpotent. Since $n\ge2$, it is nonzero. We have $1\le\dim\ker B\le\dim\ker B^2=1$, but nilpotence requires a strict increase from the first to the second kernel until reaching the full space. The first kernel has dimension one below $n$, a contradiction.

For $\lambda\ne0$, choose $\sigma^2=\lambda$ and set $B=\sigma\sum_{j=0}^{n-1}\binom{1/2}{j}(J_n/\lambda)^j$. The square-root power-series identity holds modulo $X^n$, so substitution of $J_n/\lambda$ proves $B^2=\lambda I+J_n$. This is an exact finite polynomial identity.
\end{solution}
\begin{solution}{linear:ex:commutant}
\idea Separate polynomials on individual Jordan blocks already provide $n$ independent commuting matrices.
\writeup Commutation is linear and preserved by conjugation. In Jordan coordinates, each block of size $m$ admits independent commuting matrices $I,J_m,\ldots,J_m^{m-1}$. Place each in its own diagonal block and zero elsewhere to get $n$ independent matrices. If eigenvalues are distinct, diagonal coordinates give $(\lambda_i-\lambda_j)x_{ij}=0$, so the commutant is exactly the diagonal matrices and has dimension $n$. Lagrange interpolation prescribes all their diagonal entries by one polynomial $P(\lambda_i)$, proving they are precisely $P(A)$.
\end{solution}
\begin{solution}{linear:ex:rotation}
\idea The rotation generator squares to minus identity, while a nilpotent exponential truncates.
\writeup For $K=\begin{pmatrix}0&-1\\1&0\end{pmatrix}$, even and odd powers give $e^{tK}=\cos t\,I+\sin t\,K$. Thus its matrix is $\begin{pmatrix}\cos t&-\sin t\\\sin t&\cos t\end{pmatrix}$. Since the scalar part commutes,
\[e^{t(\lambda I+J_3)}=e^{\lambda t}(I+tJ_3+t^2J_3^2/2).
\]
Its diagonal, first upper diagonal and corner are respectively $e^{\lambda t}$, $te^{\lambda t}$ and $t^2e^{\lambda t}/2$.
\end{solution}
\begin{solution}{linear:ex:minpolyjordan}
\idea Vanishing on a block means divisibility by the corresponding eigenvalue power, in every characteristic.
\writeup Expand $P(\lambda+X)=\sum_jc_jX^j$. The matrices $I,J_m,\ldots,J_m^{m-1}$ are independent, since their nonzero entries occupy different diagonals. Hence $P(\lambda I+J_m)=0$ exactly when $c_0,\ldots,c_{m-1}$ vanish, equivalently $(X-\lambda)^m\mid P$. Combining all blocks gives the least monic annihilator $\prod_\lambda(X-\lambda)^{s_\lambda}$. Its exponents equal the characteristic multiplicities precisely when each eigenvalue has one block. Using coefficients rather than derivatives avoids any factorial restriction.
\end{solution}
\begin{solution}{linear:ex:expsum}
\idea Unequal product diagonals give an exact obstruction to the exponential addition rule.
\writeup Take $A=E_{12},B=E_{21}$. Their exponentials are $I+A,I+B$, whose product is $\begin{pmatrix}2&1\\1&1\end{pmatrix}$. But $(A+B)^2=I$, so $e^{A+B}=\cosh1\,I+\sinh1\,(A+B)$ has equal diagonal entries. The matrices cannot agree, without estimating $\cosh1$ numerically. For the second example use the nonzero $A=2\pi K$ from the rotation solution; $e^A=I$ follows from sine and cosine at $2\pi$.
\end{solution}
\begin{solution}{linear:ex:symprod}
\idea Reverse each Jordan chain and transport that symmetric reversal by a congruence.
\writeup (a) Reversing the coordinate order changes the upper shift to the lower shift, so $RJ_mR=J_m^{\mathsf T}$; checking it on $e_j$ gives $e_{j+1}$, with the last vector sent to zero. (b) Write $A=PJP^{-1}$ and let $R$ reverse each block. Then $R=R^{\mathsf T}=R^{-1}$ and $RJR=J^{\mathsf T}$. Set $S=PRP^{\mathsf T}$, symmetric and invertible. Multiplication gives $S^{-1}AS=P^{-\mathsf T}J^{\mathsf T}P^{\mathsf T}=A^{\mathsf T}$. (c) The identity $AS=SA^{\mathsf T}$ makes $T=AS$ symmetric. Thus $A=TS^{-1}$ is the required product, and $S^{-1}$ is symmetric and invertible.
\end{solution}
\begin{solution}{ex:v8cCa}
\idea Projection polynomials equal one on their target primary factor and zero on the other.
\writeup Modulo $(X-1)^2$, $(X-1)^2/4$ is zero; modulo $X+1$ it is one. The complementary polynomial has the opposite residues. Since $\mu_A$ annihilates $A$, these residues prove $E_-^2=E_-$, $E_+^2=E_+$, their product zero and their sum identity. They act as the stated primary projections. Therefore $D=I-2E_-$ equals $1$ and $-1$ on the two summands. On the plus summand $N=A-I$ has square zero, and on the minus summand it is zero, proving $N^2=0$.
\end{solution}
\begin{solution}{ex:v8cCb}
\idea Squaring moves two positions at a time, separating the two parities in each chain.
\writeup The kernel increments are $4,3,2,1$, whose row lengths are $4,3,2,1$. On a length-$m$ chain, $N^2$ preserves the odd and even positions, yielding lengths $\lceil m/2\rceil$ and $\lfloor m/2\rfloor$, with an empty chain omitted. Applying the rule gives sizes $2,2$ from four, $2,1$ from three, $1,1$ from two, and one from one. Thus $N^2$ has three blocks of size two and four of size one, dimension ten and kernel dimension seven as required.
\end{solution}
\begin{solution}{ex:v8cCc}
\idea A commuting map is fixed by the image of a cyclic vector.
\writeup The vectors $J_n^je_n$, $0\le j<n$, are a basis. If $X$ commutes with $J_n$, then $XJ_n^je_n=J_n^jXe_n$, so the last column determines $X$. There are at most $n$ independent commuting matrices. The independent powers $I,J_n,\ldots,J_n^{n-1}$ already give $n$, and any last column is reached by one combination of them. Thus $X=\sum_{j=0}^{n-1}c_jJ_n^j$ uniquely, with entries $x_{il}=c_{l-i}$ for $l\ge i$ and zero below the diagonal.
\end{solution}
\begin{solution}{ex:v8cCd}
\idea Exponentiation changes a nilpotent by an invertible commuting factor, preserving every kernel level.
\writeup The finite exponential gives $e^N-I=N(I+N/2!+N^2/3!+\cdots)=NR(N)$, where the sum stops once the powers vanish. Since $R(N)=I+$ a nilpotent polynomial, a finite geometric series in $I-R(N)$ gives its inverse. It commutes with $N$. Hence $(e^N-I)^k=N^kR(N)^k$, whose kernel is exactly $\ker N^k$, since $N^kR(N)^k=R(N)^kN^k$ and the second factor on the left of $N^k$ is invertible. The Jordan rank criterion therefore gives the same nilpotent block sizes, shifted from eigenvalue zero to one.
\end{solution}

\section*{Sources and further reading}

\begin{itemize}
\item A.~Cayley, \emph{A memoir on the theory of matrices}, Philosophical Transactions of the Royal Society 148 (1858),
17--37.
\item A.-L.~Cauchy, \emph{Sur l'\'equation \`a l'aide de laquelle on d\'etermine les in\'egalit\'es s\'eculaires des
mouvements des plan\`etes}, Exercices de math\'ematiques 4 (1829). Eigenvalues of symmetric matrices are real.
\item T.~Hawkins, \emph{The Mathematics of Frobenius in Context}, Springer, 2013, and \emph{Another look at Cayley and the
theory of matrices}, Archives internationales d'histoire des sciences 27 (1977), 82--112, for the history told above.
\item F.~Brechenmacher, \emph{La controverse de 1874 entre Camille Jordan et Leopold Kronecker}, Revue d'histoire des
math\'ematiques 13 (2007), 187--257.
\item C.~Jordan, \emph{Trait\'e des substitutions et des \'equations alg\'ebriques}, Gauthier-Villars, 1870, \S\S157--159.
\item S.~Axler, \emph{Linear Algebra Done Right}, 4th ed., Springer, 2024, Chapter~8 (generalised eigenspaces and the
Jordan form without determinants).
\item R.~A.~Horn and C.~R.~Johnson, \emph{Matrix Analysis}, 2nd ed., Cambridge University Press, 2013, Chapter~3, including
the symmetric similarity of Exercise~\ref{linear:ex:symprod}.
\item G.~H.~Golub and C.~F.~Van Loan, \emph{Matrix Computations}, 4th ed., Johns Hopkins, 2013, \S7.6, on why the Jordan
form is not computed in floating point; C.~Moler and C.~Van Loan, \emph{Nineteen dubious ways to compute the exponential of a
matrix, twenty-five years later}, SIAM Review 45 (2003), 3--49.
\item G.~Frobenius, \emph{\"Uber die mit einer Matrix vertauschbaren Matrizen}, Sitzungsberichte der Preussischen Akademie
(1910), 3--15.
\item N.~J.~Fine and I.~N.~Herstein, \emph{The probability that a matrix be nilpotent}, Illinois Journal of Mathematics 2
(1958), 499--504: there are exactly $q^{n(n-1)}$ nilpotent $n\times n$ matrices over $\mathbb F_q$ .
\end{itemize}
